\documentclass[a4paper,oneside,12pt,final]{extreport}

% ============================================================
% CONFIGURATION GLOBALE (XeLaTeX)
% ============================================================
\usepackage{fontspec}
\usepackage{polyglossia}
\setmainlanguage{french}
\setotherlanguage{english}

% ============================================================
% PACKAGES DEMANDÉS + STABILITÉ DES FLOTTANTS
% ============================================================
\usepackage{graphicx}
\usepackage{float}            % autorise [H]
\usepackage[section]{placeins}% ajoute \FloatBarrier à chaque section
\usepackage{caption}
\usepackage{booktabs}
\usepackage{longtable}
\usepackage{tabularx}
\usepackage{array}
\usepackage{makecell}
\usepackage{amsmath,amsthm}
\usepackage[hidelinks]{hyperref}
\usepackage{chngcntr}
\usepackage{geometry}
\geometry{a4paper,margin=2.5cm}

% ============================================================
% NUMÉROTATION COHÉRENTE
% Figure 1.1, Figure 1.2 ; Table 1.1, Table 1.2, etc.
% ============================================================
\counterwithin{figure}{chapter}
\counterwithin{table}{chapter}
\renewcommand{\thefigure}{\thechapter.\arabic{figure}}
\renewcommand{\thetable}{\thechapter.\arabic{table}}

% ============================================================
% ENVIRONNEMENTS
% ============================================================
\newtheorem{theorem}{Theorem}[chapter]
\newtheorem{definition}{Definition}[chapter]
\newtheorem{exemple}{Example}[chapter]

% ============================================================
% AIDE COMPILATION: si une image manque, on affiche un cadre
% (évite l'échec de compilation dans les environnements incomplets)
% ============================================================
\newcommand{\safeincludegraphics}[2][]{%
  \IfFileExists{#2}{\includegraphics[#1]{#2}}{%
    \fbox{\begin{minipage}[c][4cm][c]{0.8\linewidth}\centering
      Image manquante:\\ \texttt{#2}
    \end{minipage}}%
  }%
}

\begin{document}

% ------------------------------------------------------------
% Front matter automatique (au lieu des listes écrites à la main)
% ------------------------------------------------------------
\pagenumbering{roman}


% ------------------------------------------------------------
% Pages préliminaires (version proche de ton rapport original)
% ------------------------------------------------------------
\begin{titlepage}
\centering
{\Large Projet de fin d'études}\[0.8cm]
{\bfseries\LARGE Amélioration de la gestion de maintenance de RIVA via digitalisation}\[1.2cm]
{\large Réalisé par : Ezzerki Mohamed Taha}\[0.3cm]
{\large Soutenance : 13/06/2026}\[1cm]
\vfill
{\large RIVA Industries}\
\end{titlepage}
\clearpage

\chapter*{Remerciements}
\addcontentsline{toc}{chapter}{Remerciements}
Je tiens à remercier l'ensemble des encadrants académiques et industriels pour leur accompagnement.
\clearpage

\chapter*{Résumé}
\addcontentsline{toc}{chapter}{Résumé}
Ce travail porte sur l'amélioration de la fiabilité de la machine de compactage chez RIVA Industries,
avec identification des causes d'arrêt, analyse de criticité et propositions d'amélioration.
\clearpage

\begin{otherlanguage}{english}
\chapter*{Abstract}
\addcontentsline{toc}{chapter}{Abstract}
This report focuses on improving compacting-machine reliability at RIVA Industries,
including downtime analysis, criticality assessment and maintenance improvements.
\end{otherlanguage}
\clearpage
\tableofcontents
\clearpage
\listoffigures
\clearpage
\listoftables
\clearpage

\pagenumbering{arabic}

% ============================================================
% Chapitre 1
% ============================================================
\chapter{Présentation générale de l'organisme d'accueil et cadrage du projet}

\section{Présentation de RIVA Industries}

\begin{table}[H]
\centering
\caption{Chiffres clés du groupe Meski Holding}
\label{tab:meski-cles}
\begin{tabular}{|l|l|}
\hline
\textbf{Effectif moyen} & 2000 \\
\hline
\textbf{CA consolidé} & 4,3 Milliards de MAD \\
\hline
\textbf{Capital} & 500\,000\,000 MAD \\
\hline
\textbf{Nombre de pôles d'activité} & 4 \\
\hline
\textbf{Superficie de réserve foncière} & 140 Ha \\
\hline
\end{tabular}
\end{table}

\begin{table}[H]
\centering
\caption{Fiche signalétique de RIVA Industries}
\label{tab:fiche-riva}
\begin{tabular}{|l|p{8cm}|}
\hline
\textbf{Secteur d'activité} & Sidérurgie \\
\hline
\textbf{Statut juridique} & SARL (Société à Responsabilité Limitée) \\
\hline
\textbf{Produits / Services} & Profilés et barres : Ronds à béton \\
\hline
\textbf{Maison mère ou associés} & Meski Holding (Maroc) \\
\hline
\textbf{Localisation} & Jorf Lasfar (zone de MEDZ) EL JADIDA \\
\hline
\end{tabular}
\end{table}

\begin{figure}[H]
\centering
\safeincludegraphics[width=0.9\textwidth]{media/organigramme.jpeg}
\caption{Organigramme RIVA Industries}
\label{fig:organigramme}
\end{figure}

La Figure~\ref{fig:organigramme} présente l'organigramme général.
La Table~\ref{tab:meski-cles} et la Table~\ref{tab:fiche-riva} résument les données clés.

\begin{table}[H]
\centering
\caption{Diamètre des barres de ronds à béton (RIVA1)}
\label{tab:riva1}
\begin{tabular}{|c|c|c|c|c|c|c|c|c|}
\hline
\textbf{Diamètre (mm)} & 8 & 10 & 12 & 14 & 16 & 25 & 32 & 40 \\
\hline
\end{tabular}
\end{table}

\begin{figure}[H]
\centering
\safeincludegraphics[width=0.45\textwidth]{media/image8.jpeg}
\caption{Barre à béton (produit RIVA 1)}
\label{fig:barre}
\end{figure}

\begin{figure}[H]
\centering
\safeincludegraphics[width=0.45\textwidth]{media/image9.jpeg}
\caption{Couronne rond à béton (RIVA 2)}
\label{fig:couronne}
\end{figure}

\FloatBarrier

\section{Procédure de production des couronnes}

\begin{figure}[H]
\centering
\safeincludegraphics[width=\textwidth]{media/schema_zones.jpeg}
\caption{Schéma significatif des zones de production chez RIVA 2}
\label{fig:schema-zones}
\end{figure}

\begin{figure}[H]
\centering
\safeincludegraphics[width=0.65\textwidth]{media/image12.jpeg}
\caption{Matière première -- billette}
\label{fig:billette}
\end{figure}

\begin{figure}[H]
\centering
\safeincludegraphics[width=\textwidth]{media/four_longerons.jpeg}
\caption{Four à longerons mobiles}
\label{fig:four-longerons}
\end{figure}

\begin{figure}[H]
\centering
\safeincludegraphics[width=\textwidth]{media/processus_production.png}
\caption{Processus de production}
\label{fig:processus-production}
\end{figure}

Voir la Figure~\ref{fig:schema-zones} pour les zones de production et la Figure~\ref{fig:processus-production} pour le flux global.

\FloatBarrier

% ============================================================
% Chapitre 2
% ============================================================
\chapter{Amélioration du système de maintenance}

\section{Les types de la maintenance}

\begin{figure}[H]
\centering
\safeincludegraphics[width=\textwidth]{media/types_maintenance.png}
\caption{Les types de la maintenance}
\label{fig:types-maintenance}
\end{figure}

\section{Maintenance basée sur la fiabilité}

\begin{figure}[H]
\centering
\safeincludegraphics[width=\textwidth]{media/maintenance_fiabilite.png}
\caption{Maintenance basée sur la fiabilité}
\label{fig:mbf}
\end{figure}

\section{Fonctionnement de la machine}

\begin{figure}[H]
\centering
\safeincludegraphics[width=0.6\textwidth]{media/image41.jpeg}
\caption{La compacteuse}
\label{fig:compacteuse}
\end{figure}

\begin{figure}[H]
\centering
\safeincludegraphics[width=\textwidth]{media/image42.jpeg}
\caption{Diagramme FAST}
\label{fig:fast}
\end{figure}

\section{Analyse Pareto}

\begin{figure}[H]
\centering
\safeincludegraphics[width=\textwidth]{media/pareto_2023.png}
\caption{Diagramme Pareto 2023}
\label{fig:pareto-2023}
\end{figure}

\begin{figure}[H]
\centering
\safeincludegraphics[width=\textwidth]{media/pareto_2024.png}
\caption{Diagramme Pareto 2024}
\label{fig:pareto-2024}
\end{figure}

Comparaison des résultats: voir Figure~\ref{fig:pareto-2023} et Figure~\ref{fig:pareto-2024}.

\section{Classification des composants}

\begin{figure}[H]
\centering
\safeincludegraphics[width=\textwidth]{media/image51.jpeg}
\caption{Classification des composants}
\label{fig:classification}
\end{figure}

\begin{table}[H]
\centering
\caption{Criticité/maintenance}
\label{tab:criticite-maintenance}
\begin{tabular}{|p{3cm}|p{10cm}|}
\hline
\textbf{Criticité} & \textbf{Type de maintenance} \\
\hline
\textbf{Classe AA et A} & Maintenance conditionnelle et systématique de haut niveau avec périodicité courte. \\
\hline
\textbf{Classe B} & Maintenance systématique allégée et maintenance conditionnelle selon coût. \\
\hline
\textbf{Classe C} & Travaux de base: lubrification, nettoyage, serrage. \\
\hline
\end{tabular}
\end{table}

\FloatBarrier

\chapter*{Conclusion}
\addcontentsline{toc}{chapter}{Conclusion}
Le document est maintenant prêt pour une compilation stable avec un positionnement contrôlé des figures/tableaux,
une numérotation cohérente, et des références croisées fiables.

\end{document}
